\documentclass[11pt]{article}

% ============================================================
%  Constraint-cut formula reference (quark-sector, non-custodial)
%  Each entry: the formula used for the cut, brief physical
%  intuition, and the published source.
% ============================================================

\usepackage[margin=1in]{geometry}
\usepackage{amsmath,amssymb}
\usepackage{mathtools}
\usepackage{booktabs}
\usepackage{array}
\usepackage{xcolor}
\usepackage[framemethod=default]{mdframed}
\usepackage[colorlinks=true,linkcolor=blue,citecolor=blue,urlcolor=blue]{hyperref}

\setlength{\parindent}{0pt}

\newcommand{\MKK}{M_{\mathrm{KK}}}
\newcommand{\LIR}{\Lambda_{\mathrm{IR}}}
\newcommand{\cw}{c_W}
\newcommand{\cwsq}{c_W^2}
\newcommand{\Mtwelve}{M_{12}}
\newcommand{\ImMtwelve}{\operatorname{Im}M_{12}}
\newcommand{\eps}{\varepsilon}
\newcommand{\dmk}{\Delta m_K}
\newcommand{\QVLL}{Q_1^{\mathrm{VLL}}}
\newcommand{\QVRR}{Q_1^{\mathrm{VRR}}}
\newcommand{\QLRfour}{Q_4^{\mathrm{LR}}}
\newcommand{\QLRfive}{Q_5^{\mathrm{LR}}}
\newcommand{\epsK}{\varepsilon_K}

% compact source box
\newmdenv[
  linewidth=0.5pt,
  linecolor=gray!60,
  backgroundcolor=gray!8,
  innertopmargin=4pt,
  innerbottommargin=4pt,
  innerleftmargin=6pt,
  innerrightmargin=6pt,
  skipabove=4pt,
  skipbelow=4pt
]{srcboxframe}
\newenvironment{srcbox}{\begin{srcboxframe}\small\raggedright\sloppy}{\end{srcboxframe}}

\title{\bf Warped Quark Sector Constraints}
\author{}
\date{}

\begin{document}
\maketitle
\vspace{-2.5em}

\noindent No custodial symmetry. $\MKK$ is the physical first KK gauge mass,
$\MKK = x_1\,\LIR$ with $x_1\approx2.45$. $v=246.21965$~GeV, and Wilson
coefficients are matched at $\MKK$ and QCD run to the hadronic scale.
\bigskip
% ====================================================================
\section*{Generic $\Delta F=2$}
% ====================================================================

This logic is shared by $\epsK$, $\Delta m_d$, $\Delta m_s$, and
$D^0$--$\bar D^0$.

KK-gluon exchange between quark mass eigenstates is matched onto
\begin{equation}
\mathcal{H}_{\rm eff}^{\Delta F=2}
= \sum_i C_i(\mu)\,Q_i + \sum_i \tilde C_i(\mu)\,\tilde Q_i ,
\end{equation}
with four operators: a vector-LL operator, its chirality-flipped vector-RR
partner, and two scalar left-right operators,
\begin{align}
\QVLL &= (\bar q_L^\alpha\gamma_\mu q_L^\alpha)(\bar q_L^\beta\gamma^\mu q_L^\beta), &
\QVRR &= (\bar q_R^\alpha\gamma_\mu q_R^\alpha)(\bar q_R^\beta\gamma^\mu q_R^\beta),\\
\QLRfour &= (\bar q_R^\alpha q_L^\alpha)(\bar q_L^\beta q_R^\beta), &
\QLRfive &= (\bar q_R^\alpha q_L^\beta)(\bar q_L^\beta q_R^\alpha),
\end{align}
($\alpha,\beta$ colour). The LR operators dominate; they are chirally and RG
enhanced (below).

\subsubsection*{Tree-level KK-gluon Wilson coefficients}
With $g_{L,ij}$ and $g_{R,ij}$ the off-diagonal KK-gluon couplings in the quark
mass basis (the rotated, overlap-weighted profiles times the strong coupling),
the match at $\mu=\MKK$ is
\begin{equation}
\boxed{\;
C_1^{\rm VLL}=\frac{g_{L}^2}{6\,\MKK^2},\quad
C_1^{\rm VRR}=\frac{g_{R}^2}{6\,\MKK^2},\quad
C_4^{\rm LR}=-\frac{g_{L}\,g_{R}}{\MKK^2},\quad
C_5^{\rm LR}=\frac{g_{L}\,g_{R}}{3\,\MKK^2}.
\;}
\end{equation}
Every coefficient scales as $1/\MKK^2$ (tree-level decoupling). The LR
coefficients are a product of a left and a right coupling, so both chiralities of
the flavour misalignment feed them. The colour factors ($1/6$ for VLL, $-1$ and
$1/3$ for the LR pair) are the $SU(3)$ Fierz weights.

\subsubsection*{RG running to the hadronic scale}
The coefficients run $\MKK\to2$~GeV with leading-log QCD, threshold-matched
continuously at $m_t,m_b,m_c$. VLL/VRR run multiplicatively with
$\gamma_{\rm VLL}=4$; the LR pair $\{C_4,C_5\}$ mixes under a $2\times2$
anomalous-dimension matrix that enhances the LR coefficient by a factor of several
at kaon scales. The running uses the Buras--Misiak--Urban anomalous dimensions.

\subsubsection*{Hadronic matrix elements}
For a meson $M=\bar q_1 q_2$ with decay constant $f_M$, mass $m_M$, and bag
parameters $B_1,B_4,B_5$,
\begin{align}
\langle Q_1^{\rm VLL}\rangle &= \tfrac{1}{3}\,f_M^2 m_M\,B_1, \\
\langle Q_4^{\rm LR}\rangle &= \Big(\tfrac{1}{4}\,r_\chi+\tfrac{1}{24}\Big)f_M^2 m_M\,B_4, &
\langle Q_5^{\rm LR}\rangle &= \Big(\tfrac{1}{12}\,r_\chi+\tfrac{1}{8}\Big)f_M^2 m_M\,B_5,
\end{align}
with the chiral enhancement factor
\begin{equation}
r_\chi = \left(\frac{m_M}{m_{q_1}+m_{q_2}}\right)^2 .
\end{equation}
The VRR matrix element equals the VLL one by parity.
Post-B3 audit note: this block was corrected in the June-2026 B3 audit and
confirmed by the 2026-07 full-repo audit. Older pre-B3 text used
$\tfrac23$ for $Q_1$ and had the $Q_4/Q_5$ LR coefficients swapped and doubled.

\subsubsection*{Mixing amplitude and the cut}
The new-physics mixing amplitude is the sum
\begin{equation}
\Mtwelve^{\rm NP} = \sum_i C_i\,\langle M|Q_i|\bar M\rangle
= C_1^{\rm VLL}\langle Q_1\rangle + C_1^{\rm VRR}\langle Q_1\rangle
+ C_4^{\rm LR}\langle Q_4\rangle + C_5^{\rm LR}\langle Q_5\rangle .
\end{equation}
In the repo convention the listed GeV$^3$ matrix elements are already
normalised for $M_{12}$: the usual $\langle Q_i\rangle/(2m_M)$ factor has been
absorbed, so no extra $1/(2m_M)$ is applied at this stage. The $B$ and $D$ cuts
compare $|\Mtwelve^{\rm NP}|$ to a NP budget set to half the experimental
splitting,
\begin{equation}
\boxed{\;
\begin{aligned}
\text{budget}
&=\tfrac12\,\Delta m_M^{\rm exp}
\quad(\text{or }\tfrac12|\Delta m_M^{\rm exp}-\Delta m_M^{\rm SM}|\text{ where SM is clean}),\\
\text{pass}
&\iff |\Mtwelve^{\rm NP}|\le\text{budget}.
\end{aligned}
\;}
\end{equation}

For the kaon the cut is on indirect CP violation, via the imaginary part:
\begin{equation}
\boxed{\;
\epsK^{\rm NP} = \kappa_\eps\,\frac{e^{i\phi_\eps}}{\sqrt2\,\dmk}\,
\operatorname{Im}\!\big(\Mtwelve^{\rm NP}\big),
\qquad \kappa_\eps=0.94,
\;}
\end{equation}
tested as $|\epsK^{\rm NP}|\le|\epsK^{\rm exp}-\epsK^{\rm SM}|$ (the NP budget,
$\approx6.7\times10^{-5}$). The tiny $\dmk=3.484\times10^{-15}$~GeV in the
denominator makes the kaon a strong constraint for $\Delta S=2$ CP violation, and
$\kappa_\eps=0.94$ absorbs long-distance and $\phi_\eps\neq45^\circ$ corrections.
The overall phase $e^{i\phi_\eps}$ ($\phi_\eps\approx43.5^\circ$) drops out, since
only $|\epsK^{\rm NP}|$ is compared to data.

\subsubsection*{Code evaluation}
\noindent For a sampled point the quark fit has already produced, in the down and
up sectors,
\begin{equation}
c_{Q_i},\quad c_{d_i},\quad c_{u_i},\quad
U_L^{d},\ U_R^{d},\quad U_L^{u},\ U_R^{u},
\end{equation}
together with the geometry that fixes $\MKK$. Each $\Delta F=2$ system uses the
rotations of one sector and the valence-quark pair of its meson: $(d,s)$ for
$\epsK$ and $\Delta m_K$, $(d,b)$ for $\Delta m_d$, $(s,b)$ for $\Delta m_s$, and
$(u,c)$ for $D^0$.

\smallskip
\noindent The code first gives each zero mode its squared-overlap KK-gluon
coupling $g_s^{\rm eff}F^2(c)$, so that in the gauge basis the left- and
right-handed couplings are diagonal. For the down sector,
\begin{equation}
G_L=g_s^{\rm eff}\,\mathrm{diag}\!\big(F^2(c_{Q_1}),F^2(c_{Q_2}),F^2(c_{Q_3})\big),
\qquad
G_R=g_s^{\rm eff}\,\mathrm{diag}\!\big(F^2(c_{d_1}),F^2(c_{d_2}),F^2(c_{d_3})\big),
\end{equation}
and the up sector is identical with the $c_{u_i}$. These diagonal couplings are
rotated into the quark mass basis with the matrices that diagonalise the quark mass
matrix,
\begin{equation}
\widetilde G_L=U_L^{\dagger}\,G_L\,U_L,
\qquad
\widetilde G_R=U_R^{\dagger}\,G_R\,U_R .
\end{equation}
The couplings that feed the operators are the off-diagonal entries linking the two
valence quarks,
\begin{equation}
g_L=\big(\widetilde G_L\big)_{ij},
\qquad
g_R=\big(\widetilde G_R\big)_{ij},
\end{equation}
with $(i,j)$ the generation pair of the meson. A flavour-aligned point would leave
$\widetilde G_{L,R}$ diagonal, giving $g_L=g_R=0$ and no effect; it is the anarchic
misalignment that makes them nonzero and generically complex.

\smallskip
\noindent These two numbers are inserted into the matching formulas above to set
the four Wilson coefficients
$C_1^{\rm VLL},C_1^{\rm VRR},C_4^{\rm LR},C_5^{\rm LR}$ at $\mu=\MKK$. The code then
runs them down to the hadronic scale, matching continuously at the heavy-quark
thresholds $m_t,m_b,m_c$. The vector coefficients evolve multiplicatively (the VRR
coefficient exactly like the VLL one) and the left-right pair mixes under its
$2\times2$ evolution,
\begin{equation}
C_1^{\rm VLL}(\mu)=\eta_{\rm V}\,C_1^{\rm VLL}(\MKK),
\qquad
\binom{C_4^{\rm LR}}{C_5^{\rm LR}}(\mu)
=R(\mu,\MKK)\binom{C_4^{\rm LR}}{C_5^{\rm LR}}(\MKK),
\end{equation}
with $\eta_{\rm V}$ and the matrix $R$ fixed by the QCD anomalous dimensions; $R$ is
what carries the left-right enhancement noted above.

\smallskip
\noindent The code then contracts the low-scale coefficients with the hadronic
matrix elements of the same meson (the $\langle Q_i\rangle$ assembled above from
$f_M$, $m_M$, the bag parameters, and $r_\chi$, already in the $M_{12}$
normalisation) into a single complex amplitude per point,
\begin{equation}
\Mtwelve^{\rm NP}=\sum_i C_i(\mu)\,\langle Q_i\rangle ,
\end{equation}
whose phase is inherited from the off-diagonal couplings and is generically
$\mathcal{O}(1)$.

\smallskip
\noindent The final step is the per-system comparison, using the budgets defined
above. For the $B$ and $D$ mesons the code compares the modulus to the
mass-splitting budget, $|\Mtwelve^{\rm NP}|\le\text{budget}$. For the kaon it takes
the imaginary part first, forms
$\epsK^{\rm NP}=\kappa_\eps\,\ImMtwelve^{\rm NP}/(\sqrt2\,\dmk)$, and compares
$|\epsK^{\rm NP}|$ against the indirect-CP budget
$|\epsK^{\rm exp}-\epsK^{\rm SM}|$. A point passes the $\Delta F=2$ block only when
it satisfies every system against which it is tested.

\Needspace{18\baselineskip}
\begin{srcbox}
\textbf{Source}\\[2pt]
KK-gluon flavour couplings and rotations: Csaki, Falkowski, Weiler,
arXiv:0804.1954, Sec.~2, Eqs.~(2.14)--(2.17). The $C^{1,4,5}$ tree Hamiltonian
with the $1/6$, $-1$, and $1/3$ colour factors is the displayed equation
immediately after Eq.~(2.17); Eq.~(2.18) is the chiral/RG enhancement estimate
and Eq.~(2.19) is the order-of-magnitude $C_4^K$ estimate.\\[2pt]
Running and $M_{12}$ normalisation cross-check: Blanke, Buras, Duling, Gori,
Weiler, arXiv:0809.1073, Sec.~4.2, Eqs.~(4.14)--(4.17) for the tree-level
Hamiltonian/Wilsons, Eqs.~(4.23)--(4.29) for low-scale Hamiltonian,
$M_{12}$ normalisation, and matrix elements, Eqs.~(4.36)--(4.38) for
mass-difference observables, and Eq.~(4.39) for the
$M_{12}^q=(M_{12}^q)_{\rm SM}C_{B_q}e^{2i\varphi_{B_q}}$ convention.\\[2pt]
$\epsK$ master formula
and $\kappa_\eps=0.94$: Buras, Guadagnoli, Isidori, arXiv:1002.3612, Eq.~(34) and
Eq.~(35) ($\kappa_\eps=0.94\pm0.02$); $\kappa_\eps$ introduced earlier by
Buras--Guadagnoli, arXiv:0805.3887, Eqs.~(9)--(12), lowest-order value $0.92$.
\end{srcbox}

\bigskip
% ====================================================================
\section*{$\varepsilon_K$ (kaon, CP-violating, LR-dominated)}
% ====================================================================

\begin{equation}
\boxed{\;
\big|\epsK^{\rm NP}\big|
=\left|\frac{\kappa_\eps}{\sqrt2\,\dmk}\,\operatorname{Im}\Mtwelve^{\rm NP}\right|
\;\le\; \big|\epsK^{\rm exp}-\epsK^{\rm SM}\big|,
\qquad
r_\chi=\Big(\tfrac{m_K}{m_s+m_d}\Big)^2\approx 26.
\;}
\end{equation}

\noindent The $s$--$d$ amplitude is the binding $\Delta F=2$ cut. Two things make
it strong. The kaon chiral factor $r_\chi\approx26$ multiplies the LR matrix
elements, so the LR operators ($\QLRfour,\QLRfive$) dominate $\Mtwelve^{\rm NP}$ by
orders of magnitude over the vector ones. And the anarchic KK gluon couplings carry
$\mathcal{O}(1)$ CP phases, so $\operatorname{Im}\Mtwelve^{\rm NP}$ is unsuppressed.
Inputs are the FLAG~2024 kaon bag parameters and the budget
$|\epsK^{\rm exp}-\epsK^{\rm SM}|\approx6.7\times10^{-5}$
($\epsK^{\rm exp}=2.228\times10^{-3}$ PDG, $\epsK^{\rm SM}=2.161\times10^{-3}$
Brod--Gorbahn--Stamou, who quote $2.16(18)\times10^{-3}$). The current code keeps the FLAG kaon BSM
$B_4,B_5$ inputs quoted at $3$~GeV while evolving the Wilson coefficients to
$2$~GeV; this is a known scale-matching caveat, not an additional correction
applied in the numbers above.

\begin{srcbox}
\textbf{Source}\\[2pt]
$\epsK$ formula and $\kappa_\eps$: Buras, Guadagnoli,
Isidori, arXiv:1002.3612, Eqs.~(34)--(35).\\[2pt]
Operator basis and RS matching:
Csaki--Falkowski--Weiler, arXiv:0804.1954, Sec.~2, Eqs.~(2.14)--(2.17), plus
the displayed tree Hamiltonian immediately after Eq.~(2.17).
\end{srcbox}

\bigskip
% ====================================================================
\section*{$D^0$--$\bar D^0$ (up sector)}
% ====================================================================

\begin{equation}
\boxed{\;
\big|\Mtwelve^{\rm NP}\big|_{D}\;\le\;\tfrac12\,\Delta m_D^{\rm exp},
\qquad
r_\chi=\Big(\tfrac{m_{D^0}}{m_c+m_u}\Big)^2.
\;}
\end{equation}

\noindent The same machinery applied to the $c$--$u$ mass-basis couplings: the
up-type left/right overlap matrices feed the same VLL/VRR/LR Wilson coefficients
and the same matrix-element structure with $D^0$ inputs ($f_D,m_{D^0},B_{1,4,5}$).
The $D$-meson bag inputs are weaker than the kaon and $B$ inputs in this
implementation: $B_1^D$ is a rough central value and $B_4^D=B_5^D=1$ are
internal estimates.
Because the SM short-distance $D$ mixing is swamped by long-distance
contributions, the budget is the conservative $\tfrac12\Delta m_D^{\rm exp}$ (the
whole measured splitting is allowed to be new physics, no SM subtraction). This is
the loosest of the four $\Delta F=2$ systems.

\begin{srcbox}
\textbf{Source}\\[2pt]
Same operator basis and KK-gluon matching as the master chain:
Csaki--Falkowski--Weiler, arXiv:0804.1954, Sec.~2, Eqs.~(2.14)--(2.17), plus
the displayed tree Hamiltonian immediately after Eq.~(2.17).\\[2pt]
The long-distance-dominated $D$
budget is a modelling choice.
\end{srcbox}

\bigskip
% ====================================================================
\section*{$\Delta m_d$, $\Delta m_s$ (and $\phi_s$)}
% ====================================================================

\begin{equation}
\boxed{\;
\big|\Mtwelve^{\rm NP}\big|_{B_q}\;\le\;
\max\!\Big(\tfrac12\,\Delta m_{B_q}^{\rm exp},\;
\tfrac12\big|\Delta m_{B_q}^{\rm exp}-\Delta m_{B_q}^{\rm SM}\big|\Big),
\qquad q=d,s,
\;}
\end{equation}
\begin{equation}
\boxed{\;
\phi_s = \arg\!\Big(1+\frac{\Mtwelve^{\rm NP}}{\Mtwelve^{\rm SM}}\Big).
\;}
\end{equation}

\noindent The $b$--$d$ ($B_d$) and $b$--$s$ ($B_s$) versions of the same chain.
The chiral factor is mild here ($r_\chi\approx1.6$), so the LR enhancement is much
smaller than for the kaon and the SM predictions are comparatively clean; the
budget therefore keeps an SM subtraction. The mixing phase $\phi_s$ is the argument
of the total ($\rm SM+NP$) amplitude relative to the SM, used as the $B_s$ CP-phase
observable.

\begin{srcbox}
\textbf{Source}\\[2pt]
Operator basis and KK-gluon matching:
Csaki--Falkowski--Weiler, arXiv:0804.1954, Sec.~2, Eqs.~(2.14)--(2.17), plus
the displayed tree Hamiltonian immediately after Eq.~(2.17).\\[2pt]
Running and observable cross-check: Blanke et al., arXiv:0809.1073, Sec.~4.2,
Eqs.~(4.14)--(4.17), Eqs.~(4.23)--(4.29), and Eq.~(4.39) with the following
$M_{12}^q=(M_{12}^q)_{\rm SM}C_{B_q}e^{2i\phi_{B_q}}$ convention.\\[2pt]
The $\max(\dots)$ budget and the $\phi_s$ definition
are working conventions.
\end{srcbox}

\bigskip
% ====================================================================
\section*{$\mu\to e\gamma$ (lepton flavour violation)}
% ====================================================================

\noindent The IR-brane magnetic-dipole operator and the resulting
branching ratio, with the lepton-MFV spurion $(\bar Y_N\bar Y_N^\dagger)_{12}$:
\begin{equation}
\mathcal{O}_{\rm dipole}\sim
\frac{e\,m_\mu}{2\,16\pi^2 k^2}\,(\bar Y_N\bar Y_N^\dagger)_{12}\,
F_{\mu\nu}\,\bar e_L\,\sigma^{\mu\nu}\mu ,
\end{equation}
\begin{equation}
\boxed{\;
\mathrm{BR}(\mu\to e\gamma)\simeq
4\times10^{-8}\,\big|(\bar Y_N\bar Y_N^\dagger)_{12}\big|^2
\Big(\frac{3~\mathrm{TeV}}{\MKK}\Big)^4
\;\le\; 1.5\times10^{-13},
\;}
\end{equation}
where $\bar Y_N=2k\,Y_N$ and, in the charged-lepton mass basis,
$\bar Y_N\bar Y_N^\dagger=(2k)^2\,U_{\rm PMNS}\,\mathrm{diag}(Y_{N_i}^2)\,
U_{\rm PMNS}^\dagger$.

\begin{srcbox}
\textbf{Source}\\[2pt]
Dipole operator and branching ratio: Perez \& Randall,
arXiv:0805.4652: dipole operator Eq.~(29) (full flavour structure Eq.~(28)),
branching ratio Eq.~(30) (the $4\times10^{-8}$ prefactor), spurion bound Eq.~(31)
($C\approx0.02$ in the paper-era limit). The $1.5\times10^{-13}$ bound inserted
above is the modern MEG~II 2025 90\%~C.L. limit, arXiv:2504.15711.
\end{srcbox}

\bigskip
% ====================================================================
\section*{$Z\to b\bar b$ (non-universal vertex)}
% ====================================================================

\noindent The observable is the bottom branching fraction
\begin{equation}
R_b=\frac{\Gamma(Z\to b\bar b)}{\Gamma(Z\to\text{hadrons})},
\end{equation}
(with the asymmetry $A_b$), sensitive mainly to the bottom chiral couplings
$g_L^b$ and $g_R^b$.

\smallskip
\noindent The bottom shifts (the $[2,2]$ entries of the mass-basis down-type
$Z$-coupling shifts) have two pieces. (i) The gauge-KK overlap,
a mass-basis rotation of the per-flavour gauge profile $a(c)$ relative to a fixed
EW-universal reference $a_{\rm ref}=a(0.65)$:
\begin{equation}
\delta g_L^{d}\big|_{\rm gauge}
= s_Z\,g_L^{\rm SM}\,\frac{m_Z^2}{\MKK^2}\,
U_L^\dagger\,\mathrm{diag}\!\big(a(c_{Q_i})-a_{\rm ref}\big)\,U_L ,
\end{equation}
with $g_L^{\rm SM}$ the SM left-handed chiral $Z$-coupling. (ii) The Casagrande
$m_b^2/\MKK^2$ fermion-KK admixture, summed over the relevant down-sector KK
towers:
\begin{equation}
\boxed{\;
\delta g_L^{b}\big|_{\rm admix}=\frac{m_b^2}{2\,\LIR^2}\,B_d,
\qquad
\delta g_R^{b}\big|_{\rm admix}=-\frac{m_b^2}{2\,\LIR^2}\,B_Q,
\;}
\end{equation}
\begin{equation}
\boxed{\;
B_d=\sum_i\frac{|Y_d[2,i]|^2}{|Y_d[2,2]|^2}\,B(c_{d_i},F_{d_i}),
\qquad
B_Q=\sum_i\frac{|Y_d[i,2]|^2}{|Y_d[2,2]|^2}\,B(c_{Q_i},F_{Q_i}),
\;}
\end{equation}
with the closed-form per-flavour profile
\begin{equation}
\boxed{\;
B(c,F)=\frac{1}{1-2c}\left(\frac{1}{F^2}-1+\frac{F^2}{3+2c}\right).
\;}
\end{equation}
The total shifts read by the cut are the gauge and admixture pieces added for
each chirality.

\smallskip
\noindent The IR-localised bottom fields both mix with KK gauge modes (gauge
piece) and pick up admixtures of heavy KK down-type fermions when the bottom mass
is switched on (admixture piece). The admixture pieces often dominate: the
$1/F^2$ enhancement and the sign-flipping $1/(1-2c)$ for the UV-localised light
partners make the anarchic flavour sums $B_d$ and $B_Q$ the leading effects, so
$Z\to b\bar b$ vetoes primarily through $R_b$.

\subsubsection*{Code evaluation}
\noindent For a sampled point the quark fit has already produced
\begin{equation}
c_{Q_i},\quad c_{d_i},\quad Y_d,\quad U_L^d,\quad U_R^d,\quad m_b .
\end{equation}
The code first computes one electroweak gauge-profile difference for each
left-handed quark doublet,
\begin{equation}
\alpha_{Q_i}=a(c_{Q_i})-a_{\rm ref},
\qquad
a_{\rm ref}=a(0.65).
\end{equation}
Here $a(c)$ is the numerical KK-gauge overlap sum
\begin{equation}
a(c)=\sum_n\frac{\MKK^2}{m_n^2}\,\chi_n(1)\,\Omega_n(c),
\end{equation}
so $\alpha_{Q_i}$ measures the non-universal light-$Z$ coupling of generation
$i$ relative to the reference profile.

\smallskip
\noindent The diagonal profile matrix is rotated into the down-quark mass basis,
\begin{equation}
A_L^d=(U_L^d)^\dagger
\mathrm{diag}(\alpha_{Q_1},\alpha_{Q_2},\alpha_{Q_3})U_L^d,
\end{equation}
and then multiplied by the electroweak prefactor,
\begin{equation}
\delta g_L^d\big|_{\rm gauge}
=s_Z\,g_L^{d,{\rm SM}}\left(\frac{m_Z}{\MKK}\right)^2 A_L^d,
\qquad
s_Z=-1,
\qquad
g_L^{d,{\rm SM}}=-\frac12+\frac13\sin^2\theta_W .
\end{equation}
The bottom piece is the $(3,3)$ entry, stored in the code as zero-based
index $[2,2]$,
\begin{equation}
\delta g_L^b\big|_{\rm gauge}
=\left(\delta g_L^d\big|_{\rm gauge}\right)_{33}.
\end{equation}
The code also builds the analogous right-handed shift,
\begin{equation}
\delta g_R^d\big|_{\rm gauge}
=s_Z\,g_R^{d,{\rm SM}}\left(\frac{m_Z}{\MKK}\right)^2
(U_R^d)^\dagger\mathrm{diag}\!\big(a(c_{d_i})-a_{\rm ref}\big)U_R^d,
\qquad
g_R^{d,{\rm SM}}=\frac13\sin^2\theta_W .
\end{equation}

\smallskip
\noindent If \texttt{include\_fermion\_kk\_mixing=True}, the code adds the
Casagrande fermion-KK admixture. This is where $B(c,F)$ enters: $F$ is the
IR zero-mode overlap $F(c)=f_{\rm IR}(c)$, and the per-flavour $B(c,F)$ values
are summed into
\begin{equation}
B_d=\sum_i\frac{|Y_d[2,i]|^2}{|Y_d[2,2]|^2}B(c_{d_i},F_{d_i}),
\qquad
B_Q=\sum_i\frac{|Y_d[i,2]|^2}{|Y_d[2,2]|^2}B(c_{Q_i},F_{Q_i}),
\end{equation}
\begin{equation}
\delta g_L^b\big|_{\rm admix}=\frac{m_b^2}{2\LIR^2}B_d,
\qquad
\delta g_R^b\big|_{\rm admix}=-\frac{m_b^2}{2\LIR^2}B_Q .
\end{equation}
Thus the bottom shifts passed to the observable are
\begin{equation}
\delta g_L^b
=\left(\delta g_L^d\big|_{\rm gauge}\right)_{33}
+\delta g_L^b\big|_{\rm admix},
\qquad
\delta g_R^b
=\left(\delta g_R^d\big|_{\rm gauge}\right)_{33}
+\delta g_R^b\big|_{\rm admix}.
\end{equation}
The effective $Zb\bar b$ couplings are then
\begin{equation}
\widehat g_L^b=g_L^{b,{\rm SM}}+\delta g_L^b,
\qquad
\widehat g_R^b=g_R^{b,{\rm SM}}+\delta g_R^b .
\end{equation}

\smallskip
\noindent The T010 constraint converts these couplings into
\begin{equation}
R_b^{\rm pred}=\frac{\Gamma_b^{\rm pred}}{\Gamma_{\rm had}^{\rm pred}},
\qquad
\Gamma_q\propto N_c\,\rho_q\left(|\widehat g_L^q|^2+|\widehat g_R^q|^2\right),
\end{equation}
where $\rho_q$ is the width radiator. Only the bottom coupling is shifted in
this pseudo-observable layer. The same shifted coupling gives
\begin{equation}
A_b^{\rm pred}
=\frac{|\widehat g_L^b|^2-|\widehat g_R^b|^2}
{|\widehat g_L^b|^2+|\widehat g_R^b|^2}.
\end{equation}
The measured anchors are
\begin{equation}
R_b^0=0.21629\pm0.00066,
\qquad
A_b=0.923\pm0.020 .
\end{equation}
The code forms the two pulls
\begin{equation}
p_{R_b}=\frac{R_b^{\rm pred}-R_b^{\rm exp}}{\sigma_{R_b}},
\qquad
p_{A_b}=\frac{A_b^{\rm pred}-A_b^{\rm exp}}{\sigma_{A_b}},
\end{equation}
with combined experimental plus SM-fit uncertainties
\begin{equation}
\sigma_{R_b}=\sqrt{0.00066^2+0.00013^2},
\qquad
\sigma_{A_b}=\sqrt{0.020^2+0.0001^2}.
\end{equation}
The scalar cut is
\begin{equation}
\boxed{\;
\mathrm{ratio}=\max\left(|p_{R_b}|,\ |p_{A_b}|\right)\le1 .
\;}
\end{equation}

\begin{srcbox}
\textbf{Source}\\[2pt]
Fermion-KK $m_b^2/\MKK^2$ admixture and zero-mode approximation for
$Z\bar b b$: Casagrande, Goertz, Haisch, Neubert, Pfoh, arXiv:0807.4937,
Sec.~6.4, Eq.~(170). The pseudo-observable formulae for $R_b^0$, $A_b$, and
$A_{\rm FB}^{0,b}$ are Eq.~(171), and the experimental anchors are Eq.~(173).
The all-generation $B(c,F)$ sum displayed above is the repo-local compact
implementation of the Casagrande-style fermion-admixture terms, not a
separately numbered paper equation.\\[2pt]
Custodial-protection framework for $Zb_Lb_L$ (background,
not used here): Agashe, Contino, Da Rold, Pomarol, arXiv:hep-ph/0605341.
\end{srcbox}

\subsubsection*{$\MKK$ reach}
\noindent Dominant constraint. Rigorous floor
$\MKK\approx25$--$30$~TeV in physical first-KK mass (equivalently $\sim10$--$12$~TeV
in $\LIR$ units). $Z\to b\bar b$ via $R_b$ vetoes essentially $100\%$ of evaluated
points out to $15$~TeV and $\sim89\%$ at $20$~TeV. Note the convention $\MKK=x_1\LIR$; the admixture prefactor
uses $m_b^2/(2\LIR^2)$, so the bound reads $\sim10$--$12$~TeV in Casagrande's
$M_{\rm KK}\equiv\LIR$ convention.

\bigskip
% ====================================================================
\section*{$S$, $T$, $U$ oblique parameters (minimal / non-custodial)}
% ====================================================================

\noindent The repo proxy for the minimal-RS shifts:
\begin{equation}
\boxed{\;
\Delta T_{\rm minimal}=\frac{\pi L}{2\,\cwsq}\,\frac{v^2}{\MKK^2},
\qquad
\Delta S=c_S\,\frac{v^2}{\MKK^2}\ \ (c_S=30),
\qquad
\Delta U=0,
\;}
\end{equation}
with $L=k\pi r_c\approx35$ the warped volume, $v=246.21965$~GeV, and $\cwsq=1-\sin^2\theta_W$
with $\sin^2\theta_W=0.23122$.

A predicted point $(\Delta S,\Delta T)$ is tested against
the PDG global-fit ellipse (with $U$ fixed) via the correlated Gaussian
\begin{equation}
\boxed{\;
\chi^2=(\Delta x)^\top C^{-1}(\Delta x)\le 5.991,
\qquad
\Delta x=\binom{\Delta S-0.026}{\Delta T-0.047},
\;}
\end{equation}
with $\sigma_S=0.075$, $\sigma_T=0.066$, $\rho_{ST}=0.90$, and $5.991$ the
95\%~C.L.\ threshold for two degrees of freedom.

\smallskip
\noindent New physics coupling mainly through the EW gauge self-energies shows up
as $S,T$. In the minimal bulk gauge group there is no custodial symmetry, so $T$ is
volume-enhanced by $L\approx35$ (the same log that addresses the hierarchy),
pushing the minimal-RS EW estimate to $\sim8$--$10$~TeV; $S$ is positive,
$\mathcal{O}(\text{tens})$, and not volume-enhanced.

\begin{srcbox}
\textbf{Source}\\[2pt]
Prediction side: repo proxy \texttt{quarkConstraints.oblique\_stu.rs\_minimal\_oblique\_proxy}; no full
warped EW fit is performed. The $c_S=30$ coefficient is taken from the PDG~2025
Electroweak review's warped-extra-dimension estimate
$S\simeq30\,v^2/M_{\rm KK}^2$, which cites Carena et al.,
Nucl.~Phys.~B~759~(2006)~202, arXiv:hep-ph/0607106. Agashe, Delgado, May,
Sundrum, arXiv:hep-ph/0308036, Sec.~5.1 documents the
custodial/volume-enhancement context (Eqs.~(5.5)--(5.6)).\\[2pt]
Fit ellipse: PDG~2025 Standard Model review, Table~10.8
($S=0.026\pm0.075$, $T=0.047\pm0.066$, $\rho=0.90$, $U$ fixed). The
95\% two-parameter threshold is $\chi^2=5.991$.
\end{srcbox}

\subsubsection*{$\MKK$ reach}
$\MKK\approx5$--$7$~TeV kills $100\%$ of
points at the $5$~TeV tile (all $r$) and relaxes by $7$~TeV. This is a proxy (the
prediction side is the parametric estimate above; the data $\chi^2$ is rigorous),
so it does not contribute to the rigorous floor; it is part of the $\sim7$~TeV proxy
reach but subdominant to the rigorous $Z\to b\bar b$ ($\sim25$--$30$~TeV).
\bigskip
% ====================================================================
\section*{Direct collider bounds on KK masses}
% ====================================================================

\noindent No formula: these are numeric mass-exclusion edges from
ATLAS/CMS searches, applied as $\MKK > m_{\rm excl}$. The active values are:

\smallskip
\begin{center}
\small
\renewcommand{\arraystretch}{1.15}
\begin{tabular}{
  >{\raggedright\arraybackslash}p{0.35\textwidth}
  >{\raggedright\arraybackslash}p{0.16\textwidth}
  >{\raggedright\arraybackslash}p{0.35\textwidth}}
\toprule
Process & Excluded mass & Cited search \\
\midrule
KK-gluon $\to t\bar t$            & up to $5.5$~TeV  & CMS-B2G-25-009, arXiv:2603.23454 \\
$T_{5/3}$ exotic top partner (pair) & $\gtrsim1.46$~TeV & PDG q009 / ATLAS arXiv:2212.05263 \\
charge-$2/3$ top partner $T$ (pair) & $\gtrsim1.70$~TeV & PDG Q009TPP / ATLAS arXiv:2401.17165 \\
charge-$1/3$ bottom partner $B$ (pair) & $\gtrsim1.57$~TeV & PDG~2025 $b'$ / CMS PRD~110~(2024)~052004 \\
dilepton resonance ($Z'_{\rm SSM}$) & $\gtrsim5.15$~TeV & CMS JHEP~07~(2021)~208 \\
charged-current resonance ($W_{\rm KK}/W'$) & $\gtrsim6.0$~TeV & PDG~2025 $W'$ / ATLAS arXiv:1906.05609 \\
spin-2 RS graviton ($k/\bar M_{\rm Pl}=0.1$) & $\gtrsim4.8$~TeV & PDG~2025 extra dimensions / CMS diphoton \\
singlet VLQ $T$ (non-custodial, pair) & $\gtrsim1.36$~TeV & ATLAS 2024 (arXiv:2401.17165) \\
Drell--Yan contact operator ($\Lambda_{LL}^+$) & $\gtrsim35.8$~TeV & PDG~2025 compositeness \\
VLQ doublet $(T,B)$ (pair)       & $\gtrsim1.37$~TeV & PDG~2025 (T,B)-doublet \\
longitudinal VBS ($W_LW_L$)      & $\sigma$-limit (no mass edge) & ATLAS PRL~135~(2025)~111802 \\
diboson spin-1 ($W'/V'$)         & $\gtrsim4.4$--$4.8$~TeV & CMS-B2G-20-009 Table~2 \\
diphoton resonance ($G^*$)       & $\gtrsim4.8$~TeV & PDG~2025 extra dimensions / CMS arXiv:2405.09320 \\
four-top ($Z'$)                  & up to $0.85$~TeV & CMS-B2G-25-005 / arXiv:2604.14058 \\
\bottomrule
\end{tabular}
\end{center}

\smallskip
\noindent These are direct-search exclusion edges reinterpreted as lower bounds on
the relevant KK/partner mass. The Drell--Yan contact operator and longitudinal VBS
rows are not single mass edges (a contact-operator scale and a fiducial
cross-section limit, respectively) but are listed for completeness.

\begin{srcbox}
\textbf{Source}\\[2pt]
These rows are not equation citations. Each number is the scalar edge tracked
in the local collider catalog entries
\texttt{flavor\_catalog/processes/collider\_rs/CR001}--\texttt{CR014}, with
source snapshots under \texttt{flavor\_catalog/references/CR*/}. Rows marked
``contact operator'' or ``$\sigma$-limit'' are included for context and are not
single $\MKK$ cuts.
\end{srcbox}

\subsubsection*{$\MKK$ reach}
\noindent Floor $\MKK\approx5$~TeV, set by the KK-gluon
$\to t\bar t$ search, which vetoes $100\%$ of points out to $5$~TeV. Only the lepton-free searches enter the quark-only
scan: the lepton-involving dilepton, $W'\to\ell\nu$, Drell--Yan contact, VBS, and
four-top two-lepton searches are excluded, so their exclusion edges set no floor
here. 

\bigskip
% ====================================================================
\section*{$\MKK$ floors}
% ====================================================================

\noindent The floors are the physical first-KK mass $\MKK=x_1\LIR$
($x_1\approx2.45$) at or above which each process stops excluding points.
\smallskip
\begin{center}
\small
\begin{tabular}{
  >{\raggedright\arraybackslash}p{0.39\textwidth}
  >{\raggedright\arraybackslash}p{0.17\textwidth}
  >{\raggedright\arraybackslash}p{0.31\textwidth}}
\toprule
Process & Tag & $\MKK$ floor (physical) \\
\midrule
$\varepsilon_K$        & rigorous & $\sim2$--$3$~TeV  \\
$\Delta m_K$           & rigorous & no floor (never vetoes) \\
$D^0$--$\bar D^0$      & rigorous (CPV partner proxy) & no floor (never vetoes) \\
$\Delta m_d$           & rigorous & $\sim1$--$2$~TeV  \\
$\Delta m_s$           & rigorous & $\sim1$--$2$~TeV  \\
$\phi_s$               & rigorous & $\sim2$--$3$~TeV  \\
combined rigorous $\Delta F=2$ & rigorous & $\sim1$--$3$~TeV \\
$\mu\to e\gamma$       & n/a & n/a (in progress) \\
$Z\to b\bar b$ (\textbf{dominant}) & rigorous & $\sim25$--$30$~TeV (phys.) \\
\quad ($Z\to b\bar b$ in $\LIR$ units) & rigorous & $\sim10$--$12$~TeV ($\LIR$) \\
$S,T,U$ oblique        & proxy & $\sim5$--$7$~TeV  \\
collider (KK-gluon $\to t\bar t$) & proxy & $\sim5$~TeV \\
collider (diboson/diphoton) & proxy & $\sim3$~TeV \\
collider (VLQ pairs)   & proxy & $\sim1$~TeV \\
\bottomrule
\end{tabular}
\end{center}


\end{document}
